\documentclass[10pt,a4paper]{amsart}
\usepackage[margin=19mm]{geometry}
\usepackage{amsmath,amssymb,mathtools}
\usepackage[scr=rsfs,cal=euler]{mathalfa}
\usepackage{xfrac}
\usepackage[unicode,hidelinks,bookmarks=false]{hyperref}
\newcommand{\ie}{i.e.\ }
\newcommand{\cf}{cf.\ }
\newcommand{\resp}{respectively\ }
\newcommand{\Subsection}{\subsection}
\providecommand{\nobreakdash}{\nobreak\hbox{-}}
\setlength{\emergencystretch}{2em}
\multlinegap0pt
\usepackage{bm}
\usepackage{bbm}
\usepackage{dsfont}
\usepackage{xspace}
\usepackage{cancel}
\usepackage{mathtools}
\usepackage{amsthm}
\makeatletter
\@ifundefined{Theorem}{\newtheorem{Theorem}{Theorem}}{}
\@ifundefined{Lemma}{\newtheorem{Lemma}[Theorem]{Lemma}}{}
\@ifundefined{Proposition}{\newtheorem{Proposition}[Theorem]{Proposition}}{}
\makeatother
\newcommand{\cR}{\ensuremath{\mathbbm{R}}}
\newcommand{\eps}{\varepsilon }
\newcommand{\reg}{\mathrm{reg}}
\newcommand{\tmmathbf}[1]{\ensuremath{\boldsymbol{#1}}}
\newcommand{\vf}{f}
\title[Piecewise Smooth Dynamical Systems Regularized by Convolutio]{Piecewise Smooth Dynamical Systems Regularized by Convolution}
\author{Claudio A. Buzzi, Daniel Panazzolo, Paulo R. da Silva}
\date{Source: arXiv:2601.00697v2 (CC BY 4.0)}
\begin{document}
\maketitle

\noindent\emph{Source and licence.} Piecewise Smooth Dynamical Systems Regularized by Convolution. Version arXiv:2601.00697v2 (CC BY 4.0), distributed under the Creative Commons Attribution 4.0 International licence, as recorded in the arXiv licence metadata for this exact version and shown on the abstract page https://arxiv.org/abs/2601.00697v2. Full rights notice: licenses/arxiv-2601.00697-CC-BY-4.0.txt.\par\smallskip

\noindent\textit{Editorial scope.} Counted unit: the Proposition whose source label is \texttt{prop-smoothing-one-comp} in the article (source0.tex lines 1558--1562, source label \texttt{prop-smoothing-one-comp}), delivered together with the article's own complete original proof of it (source0.tex lines 1563--1660, one \texttt{proof} environment of 98 lines). No mathematics is written, repaired or re-derived: statement and proof are the article's own text line for line. Required context retained uncounted: expr-f\_k (lines 647--649); Lemma-partial-conv (lines 1137--1139); def-X (lines 1536--1538). Every definition, display and label of those lines is retained unchanged. Omitted from this package and not redistributed: 1--646, 650--1136, 1140--1535, 1539--1557, 1661--2371. Nothing inside a retained excerpt is omitted or altered: every retained body is delivered exactly as the article's lines are written, so no omission receipt of any kind accompanies this package. Citations: the retained excerpts carry no citation command at all, so no bibliography entry of the article is transcribed and no reference is redistributed. Reference adaptation: none was needed and none was made; every reference inside the delivered unit resolves inside it, so no unresolved reference or citation is produced. Printed slips kept literal: no printed word, symbol, hypothesis or number of the counted statement or of its proof was altered. Layout and class: the article's own class is \texttt{article} with packages including babel, amsmath, amssymb, graphicx, enumerate, bbm, accents, caption, amsthm, stmaryrd, float, appendix, array, dsfont; the wrapper is the lane's pinned \texttt{amsart} editorial wrapper at 10pt on A4, which restates the theorem-like environments the retained text needs and adds the article's own macro definitions that the retained text needs (\texttt{\textbackslash cR}, \texttt{\textbackslash eps}, \texttt{\textbackslash reg}, \texttt{\textbackslash tmmathbf}, \texttt{\textbackslash vf}), with extra wrapper packages (bm, bbm, dsfont, xspace, cancel, mathtools). The delivered numbering is the unit's own. Assets: no retained span contains a figure environment, a graphics-inclusion command, a PDF-page inclusion, a table or a TikZ picture; no image file is copied into this package, the package declares no asset dependency and no image file is redistributed with it. Licence: version arXiv:2601.00697v2 of this article is distributed under the Creative Commons Attribution 4.0 International licence, as recorded in the arXiv licence metadata for this exact version and shown on the abstract page https://arxiv.org/abs/2601.00697v2. A full rights notice is delivered at licenses/arxiv-2601.00697-CC-BY-4.0.txt. Characters: every retained excerpt is pure ASCII, so the delivered engine, XeLaTeX with its default Latin Modern text font, reports no missing character.
\begin{equation}\label{expr-f_k}
m_\eps \ast \vf(x)= \int_{\cR^n} f (x - \eps t)\, m (t) d t
\end{equation}

\begin{Lemma}\label{Lemma-partial-conv}
The function $f^\reg\circ B_F$ is smooth.
\end{Lemma}

\begin{equation}\label{def-X}
 X = \tmmathbf{1}_{\{ x_1 > 0 \}} X^+ + \tmmathbf{1}_{\{ x_1 < 0 \}} X^- 
\end{equation}

\begin{Proposition}\label{prop-smoothing-one-comp}
The vector field 
$$\mathcal{X} = r\cdot (\varphi^\star X^\reg)$$ 
defines an element of $\mathfrak{X }^{\infty}(\tilde N,\tilde M,\emptyset,\tilde \pi)$.  In particular, $\mathcal{X}$ is smooth.
\end{Proposition}

\begin{proof}
We firstly study the expression of the regularized vector field $\mathcal{X}$ in the initial coordinates. 
By hypothesis, the coefficients of the original vector field $X$ in (\ref{def-X}) can be written as
\begin{equation}\label{coeff-fkX}
\vf_k = \tmmathbf{1}_{\{ x_1 > 0 \}} \vf_{k }^+ + \tmmathbf{1}_{\{ x_1 < 0 \}} \vf_{k }^- 
\end{equation}
where both $\vf_k^+$ and $\vf_k^-$ are smooth functions on $\cR^n$. 
By Fubini's Theorem, the integral expression (\ref{expr-f_k}) for $f_k$ can be rewritten as
\begin{equation}\label{Fubini-f_k}
  f_k \left( x, \underline{x}, \eps  \right) = \int_{\cR^{n - 1}}
  \left( \int_{\cR}  \vf_k \left( x_1- \eps \tau, \underline{x} -
  \eps  \underline{\tau} \right) m \left( \tau, \underline{\tau} \right)
  d \tau \right) d \underline{\tau}
\end{equation}
where we write $x = (x_1,\underline{x})$ and $t = (\tau,\underline{\tau}) \in \cR \times \cR^{n-1}$. Using the above decomposition of $\vf_k$, the innermost integral can be further decomposed as 
{\small \begin{align}\label{equation-integexpr}
\int_{\cR} 
   \vf_k \left( x_1- \eps \tau, \underline{x} -
  \eps  \underline{\tau} \right) m \left( \tau, \underline{\tau} \right) d\tau
   &= \int_{\{ x_1- \eps \tau \ge 0 \}} 
        \vf_k^+ \left( x_1- \eps \tau, \underline{x} -
  \eps  \underline{\tau} \right) m \left( \tau, \underline{\tau} \right)  \\
   &\;\;+ \int_{\{ x_1- \eps \tau \le 0 \}} 
        \vf_k^- \left( x_1- \eps \tau, \underline{x} -
  \eps  \underline{\tau} \right) m \left( \tau, \underline{\tau} \right) \notag
\end{align}}
which obviously has a discontinuous limit as $\eps \rightarrow 0$. 

Let us now compute the pull-back of $X^\reg$ under the blowing-up through the coordinates of the directional charts.  Firstly, the $\eps$-directional chart is given by the coordinate change $x_1 = \eps  y$, with a new variable $y\in \cR$.  In this chart, the vector field assumes the form 
\[ 
 \frac{1}{\eps } F_1 \left( y, \underline{x}, \eps 
   \right) \frac{\partial}{\partial y} + \sum_{k \ge  2} F_k \left( y,
   \underline{x}, \eps  \right) \frac{\partial}{\partial x_k} \]
where each function $F_k = f_k \circ \varphi$ is given $F_k \left( y, \underline{x}, \eps  \right) : = f_k \left(
\eps  y, \underline{x}, \eps  \right)$. Therefore, since the divisor $D$ is expressed in this chart by the equation $\{\eps = 0\}$, we conclude that ${\mathcal{X}}$ 
is locally generated by the vector field
\begin{equation}\label{expression-calX}
  F_1 \left( y, \underline{x}, \eps 
   \right) \frac{\partial}{\partial y} + \eps \left(\sum_{k \ge  2} F_k \left( y,
   \underline{x}, \eps  \right) \frac{\partial}{\partial x_k}\right)
\end{equation}
As a consequence, it suffices to prove that each coefficient $F_k$ extends smoothly to $\{\eps = 0\}$. 
By substituting $x_1=\eps  y$, the integral (\ref{equation-integexpr}) assumes the form
\begin{align}\label{equation-integexprF}
\int_{\cR} 
   \vf_k\!\left( \eps  (y - \tau), \, \underline{x} - \eps  \underline{\tau} \right) 
   m \!\left( \tau, \underline{\tau} \right) d\tau
   &= \int_{\{ \tau \le  y \}} 
        \vf_k^+ \!\left( \eps  (y - \tau), \, \underline{x} - \eps  \underline{\tau} \right) 
        m \!\left( \tau, \underline{\tau} \right) d\tau  \\
   &\; + \int_{\{ \tau \ge  y \}} 
        \vf_k^- \!\left( \eps  (y - \tau), \, \underline{x} - \eps  \underline{\tau} \right) 
        m \!\left( \tau, \underline{\tau} \right) d\tau \notag
\end{align}
And, by substituting $x_1=\eps  y$ into the the integral (\ref{Fubini-f_k}), we conclude that $F_k$ is given by the integral
$$
F_k(y,\underline{x},\eps) = \int_{\cR^{n-1}} \; \int_{\cR} 
   \vf_k\!\left( \eps  (y - \tau), \, \underline{x} - \eps  \underline{\tau} \right) 
   m \!\left( \tau, \underline{\tau} \right) d\tau\; d\underline{\tau}
$$
which is precisely the integral studied in Lemma \ref{Lemma-partial-conv} in the particular case where $I=\{1\}$. Therefore, $F_k$ is a smooth function. 

We now consider the $\pm x_1$-directional charts, where the blow-up takes the form $x_1 = \pm z, \eps= z \rho$, with new variables $z,\rho\in \cR_{\ge 0}$. Using these coordinates, the vector field $\mathcal{X}$ (outside the divisor $D$) is given by
\[
 \mathcal{X}= \frac{1}{z} G_{1,\pm} \left( z, \underline{x}, \rho
   \right) \left(z\frac{\partial}{\partial z}-\frac{\partial}{\partial \rho}\right) + \sum_{k \ge  2} G_{k,\pm} \left( z,
   \underline{x}, \rho \right) \frac{\partial}{\partial x_k} 
\]
where $G_{k,\pm}=f_k \circ \varphi$, i.e.\ $G_{k,\pm}( z,\underline{x}, \rho) = f_k(\pm z,\underline{x},z \rho)$.  
In this chart, the divisor has equation $D=\{ z = 0\}$ and therefore ${\mathcal{X}}$ is equivalent to 
$$
G_{1,\pm} \left( z, \underline{x}, \rho
   \right) \left(z\frac{\partial}{\partial z}-\frac{\partial}{\partial \rho}\right) + z \left( \sum_{k \ge  2} G_{k,\pm} \left( z,
   \underline{x}, \rho \right) \frac{\partial}{\partial x_k}  \right)
$$
We claim that each coefficient $G_{k,\pm}$ is smooth.  To simplify the notation, we make the computations for $G_{k,+}$, the negative case being similar.

Referring again to the integral in (\ref{equation-integexpr}), the substitution $x_1 = z, \eps= z \rho$ yields
\begin{align*}
\int_{\cR} 
   \vf_k\!\left( z(1-\rho), \, \underline{x} - z \rho \underline{\tau} \right) 
   m \!\left( \tau, \underline{\tau} \right) d\tau
   &= \int_{{\{ \rho \tau \le 1 \}}} 
        \vf_k^+ \!\left( z(1-\rho \tau), \, \underline{x} - z \rho \underline{\tau} \right) 
        m \!\left( \tau, \underline{\tau} \right) d\tau  \\
   &\; + \int_{{\{ \rho\tau  \ge 1 \}}} 
        \vf_k^- \!\left( z(1-\rho \tau), \, \underline{x} - z \rho \underline{\tau} \right) 
        m \!\left( \tau, \underline{\tau} \right) d\tau
\end{align*}
We claim that this is smooth function of the variables $z,\rho,\underline{x}$ and $\underline{\tau}$. To see this, the crucial fact to observe is that for $\rho$ sufficiently small, the second integral in the right-hand side {\em vanish identically}. Indeed, since the mollifier $m$ has compact support, there exists an $\rho_0 > 0$ such that $m(\tau,\underline{\tau})$ vanishes identically for $\tau \ge \frac{1}{\rho_0}$. Therefore, for all $\rho \le \rho_0$, the above integral assumes the simple form 
$$
\int_{{\cR}} 
        \vf_k^+ \!\left( z(1-\rho \tau), \, \underline{x} - z \rho \underline{\tau} \right) 
        m \!\left( \tau, \underline{\tau} \right) d\tau
$$
which is a smooth function since $\vf_k^+$ is smooth. 
A further integration with respect to $\underline{\tau}$ shows that $G_{k,+}$ is smooth. This concludes the proof. 
\end{proof}

\end{document}
